% ================================================================
% frame1-tex.tex
% The Numerical Framework — Al-Hindara Mother Paper
% Author: Eng. Abdulla M. N. Al-Jabri — Sana'a, Yemen
% Date: October 2026
% Status: Conceptual and exploratory mathematical framework
% ================================================================

\documentclass[11pt,a4paper]{article}

% ------------------------------------------------
% Packages
% ------------------------------------------------
\usepackage[utf8]{inputenc}
\usepackage[T1]{fontenc}
\usepackage{amsmath,amssymb,amsthm}
\usepackage{geometry}
\usepackage{graphicx}
\usepackage{hyperref}
\usepackage{booktabs}
\usepackage{xcolor}
\usepackage{float}
\usepackage{fancyhdr}

% ------------------------------------------------
% Page layout and metadata
% ------------------------------------------------
\geometry{margin=1in}

\hypersetup{
colorlinks=true,
linkcolor=blue,
citecolor=blue,
urlcolor=cyan,
pdftitle={The Numerical Framework: Riemann Zeta Zeros and the Proposed Cutoff 21},
pdfauthor={Eng. Abdulla M. N. Al-Jabri},
pdfsubject={DOI: 10.5281/zenodo.23540283}
}

% ------------------------------------------------
% Header and Footer Customization
% ------------------------------------------------
\pagestyle{fancy}
\fancyhf{}
\fancyhead[L]{\small \textbf{Eng. Abdulla M. N. Al-Jabri} -- Frame 1}
\fancyhead[R]{\small Al-Hindara Mother Paper}
\fancyfoot[C]{\thepage}
\fancyfoot[L]{\small\textit{Al-Hindara Numerical Framework}}
\fancyfoot[R]{\small \texttt{jabri62018@gmail.com}}
\renewcommand{\headrulewidth}{0.4pt}
\renewcommand{\footrulewidth}{0.4pt}
\setlength{\headheight}{14pt}

% ------------------------------------------------
% Title
% ------------------------------------------------
\title{
\textbf{The Numerical Framework:}\\
\Large A Damping Model, the Riemann Zeta Zeros,
and the Proposed Integer Cutoff \(x_p=21\)\\[0.5em]
\large A Conceptual and Exploratory Study within
the Al-Hindara Framework
}

\author{
\textbf{Eng. Abdulla Mohammed Nasser Al-Jabri}\\
\small Independent Researcher -- Sana'a, Yemen\\
\small\texttt{jabri62018@gmail.com}\\
\small ORCID:
\href{https://orcid.org/0009-0003-3319-3822}
{0009-0003-3319-3822}
}

\date{
October 2026 \\[1em]
\small \textbf{DOI:} \href{https://doi.org/10.5281/zenodo.23540283}{10.5281/zenodo.23540283} \\
\small \textbf{Zenodo Record:} \url{https://zenodo.org/records/23540283}
}

\begin{document}

\maketitle
\thispagestyle{fancy}

% ================================================================
% ABSTRACT
% ================================================================
\begin{abstract}
\noindent
This paper presents a conceptual numerical framework within
the Al-Hindara project. It introduces the normalization proposal
\(Z+C+A=1\) and studies a candidate damping function involving
the integer \(x_p=21\), defined as the floor of the ordinate of
the second positive nontrivial zero of the Riemann zeta function.

The proposed function is
\[
Z(x)=x^{-5}\ln(x)\sin\left(\frac{2\pi}{x}\right)e^{-x/x_p},
\qquad x>1.
\]
The choice \(x_p=21\) is a modeling decision motivated by the
numerical value \(\gamma_2\approx21.022039\). It is not a
consequence of a theorem about the zeta function.

For the accumulation function
\[
A(x)=\int_1^x Z(t)\,dt,
\]
sign analysis shows that \(A\) decreases on \(1<x<2\) and
increases on \(x>2\). Consequently, the proposed function does
not have a maximum at \(x=42\) on the interval \(1<x\leq200\).
This distinction between a numerical proposal and a verified
mathematical property is made explicit.

The paper also discusses the quantity
\(K=2\pi e^{-1/42}\) and its proposed role in the model.
No asymptotic connection between this quantity and the
Riemann zeta function is established here.

A companion study of tree crown ratios is discussed as an
exploratory empirical investigation. Its small sample does not
establish a causal or universal biological relationship.
The paper therefore separates exact mathematical results,
reported empirical observations, and conjectural modeling
proposals.
\vspace{1em}

\noindent\textbf{Keywords:}
Riemann zeta function, nontrivial zeros, damping function,
integer cutoff, Al-Hindara, numerical modeling,
reciprocal Silver Ratio, exploratory analysis.
\end{abstract}

% ================================================================
% 1. INTRODUCTION
% ================================================================
\section{Introduction}

The Al-Hindara project proposes the conceptual normalization
relation
\begin{equation}
Z+C+A=1.
\label{eq:normalization}
\end{equation}

In this paper, Equation~\eqref{eq:normalization} is treated
as a conceptual framework rather than an established
mathematical identity. A rigorous quantitative interpretation
would require explicit definitions of \(Z\), \(C\), and \(A\),
their domains, their units where applicable, and the
normalization procedure that makes their sum equal to one.

The present paper investigates a candidate numerical function
with an exponential damping term. The selected damping
parameter is
\[
x_p=21.
\]
This value is motivated by the integer part of the second
positive nontrivial zero ordinate of the Riemann zeta function.

The objectives are:

\begin{enumerate}
\item Define the numerical cutoff \(x_p\) precisely.
\item Introduce and analyze the proposed damping function.
\item Examine the mathematical behavior of its accumulation
function.
\item Distinguish numerical observations from proved results.
\item Describe the proposed relationship with the companion
tree-shadow study without asserting an unproved causal link.
\end{enumerate}

The paper does not claim to prove the Riemann Hypothesis,
derive the proposed function from the zeta function, or
establish a universal law connecting number theory and
biological morphology.

% ================================================================
% 2. THE RIEMANN ZERO AND THE CUTOFF
% ================================================================
\section{The Riemann Zero and the Proposed Integer Cutoff}

Let \(\zeta(s)\) denote the Riemann zeta function. The first
positive ordinates of its nontrivial zeros on the critical
line are conventionally denoted by \(\gamma_1,\gamma_2,\ldots\),
where a zero is written as
\[
s=\frac12+i\gamma.
\]

The first two positive ordinates are approximately
\[
\gamma_1\approx14.134725,
\qquad
\gamma_2\approx21.022039.
\]
See Edwards~\cite{edwards1974} for background on the
Riemann zeta function and its zeros.

We define the candidate integer cutoff by
\begin{equation}
x_p:=\lfloor\gamma_2\rfloor=21.
\label{eq:xp}
\end{equation}

Within the Al-Hindara project, the integer \(21\) is referred
to as the \textit{Yemeni Constant}. This is a project-specific
name for a model parameter, not a recognized mathematical
constant established by this paper.

The choice in Equation~\eqref{eq:xp} is a modeling convention.
Taking the floor of \(\gamma_2\) gives \(21\), but this operation
alone does not imply that the resulting integer has a special
analytic role in the distribution of zeta zeros.

Any such role would require a separate mathematical argument
or reproducible numerical evidence against an appropriate
comparison model.

% ================================================================
% 3. THE DAMPING FUNCTION
% ================================================================
\section{Definition of the Damping Function}

For \(x>1\), define
\begin{equation}
Z(x)=x^{-5}\ln(x)
\sin\left(\frac{2\pi}{x}\right)e^{-x/x_p},
\qquad x_p=21.
\label{eq:field}
\end{equation}

Equivalently,
\begin{equation}
Z(x)=x^{-5}\ln(x)
\sin\left(\frac{2\pi}{x}\right)e^{-x/21}.
\label{eq:field_expanded}
\end{equation}

The expression contains four main factors:

\begin{itemize}
\item \(x^{-5}\): an algebraic factor that decreases with \(x\).
\item \(\ln(x)\): a logarithmic factor.
\item \(\sin(2\pi/x)\): an oscillatory factor whose argument
depends on the reciprocal of \(x\).
\item \(e^{-x/21}\): an exponential damping factor.
\end{itemize}

The sine term should not be described as a sinusoid with a
constant period in \(x\). Its argument is \(2\pi/x\), so the
oscillation is not periodic with a fixed period on the
\(x\)-axis.

Equation~\eqref{eq:field} defines a mathematical function.
Its interpretation as a physical field, a biological model,
or a quantity derived from the Riemann zeta function has not
been established.

% ================================================================
% 4. SIGN AND ASYMPTOTIC ANALYSIS
% ================================================================
\section{Sign and Asymptotic Behavior}

For \(1<x<2\), the angle \(2\pi/x\) lies between \(\pi\) and
\(2\pi\). Therefore,
\[
\sin\left(\frac{2\pi}{x}\right)<0.
\]
Since \(x^{-5}>0\), \(\ln(x)>0\), and \(e^{-x/21}>0\) for \(x>1\),
it follows that
\begin{equation}
Z(x)<0,\qquad 1<x<2.
\label{eq:negative}
\end{equation}

At \(x=2\),
\begin{equation}
Z(2)=0.
\label{eq:zero2}
\end{equation}

For \(x>2\), the angle \(2\pi/x\) lies between \(0\) and \(\pi\),
so
\begin{equation}
Z(x)>0,\qquad x>2.
\label{eq:positive}
\end{equation}

These sign properties follow directly from the definition
and do not depend on numerical simulation.

As \(x\to\infty\), the small-angle approximation gives
\[
\sin\left(\frac{2\pi}{x}\right)
\sim\frac{2\pi}{x}.
\]
Consequently,
\begin{align}
Z(x)
&\sim x^{-5}\ln(x)\frac{2\pi}{x}e^{-x/21}\\
&=2\pi x^{-6}\ln(x)e^{-x/21}.
\end{align}

Thus,
\begin{equation}
Z(x)\sim
2\pi x^{-6}\ln(x)e^{-x/21},
\qquad x\to\infty.
\label{eq:asymptotic}
\end{equation}

This is an asymptotic statement about the defined function
\(Z(x)\). It is not an asymptotic formula for the Riemann
zeta function or its zeros.

% ================================================================
% 5. ACCUMULATION FUNCTION
% ================================================================
\section{The Accumulation Function}

Define the continuous accumulation function
\begin{equation}
A(x)=\int_1^x Z(t)\,dt,\qquad x>1.
\label{eq:accumulation}
\end{equation}

By the Fundamental Theorem of Calculus,
\begin{equation}
A'(x)=Z(x).
\label{eq:accumulation_derivative}
\end{equation}

Combining Equation~\eqref{eq:accumulation_derivative} with
the sign analysis in Section~4 gives
\begin{align}
A'(x)&<0,\qquad 1<x<2,\\
A'(2)&=0,\\
A'(x)&>0,\qquad x>2.
\end{align}

Therefore, \(A(x)\) decreases on \((1,2)\) and increases on
\((2,\infty)\). The point \(x=2\) is a minimum, not a maximum.

In particular, on the interval \(1<x\leq200\), the maximum
of \(A(x)\) occurs at the right endpoint \(x=200\), not at
\(x=42\).

Hence the claim that the accumulation function defined by
Equation~\eqref{eq:accumulation} peaks at \(2x_p=42\) is
not supported by the stated equations.

A different accumulation measure might have an extremum
near \(42\), but it would need to be defined explicitly and
analyzed independently. It would not be valid to attribute
such an extremum to Equation~\eqref{eq:accumulation} without
a corresponding derivation.

% ================================================================
% 6. A CANDIDATE NUMERICAL COEFFICIENT
% ================================================================
\section{A Candidate Numerical Coefficient}

The project also considers the quantity
\begin{equation}
K=2\pi e^{-1/42}.
\label{eq:K}
\end{equation}

Numerically,
\begin{equation}
K\approx6.1353.
\end{equation}

Equation~\eqref{eq:K} defines a valid numerical expression.
However, the expression does not follow from the definition
of \(Z(x)\) or from the accumulation function in
Equation~\eqref{eq:accumulation}.

The use of \(42=2x_p\) in Equation~\eqref{eq:K} is a modeling
choice. No result in the present paper establishes \(K\) as
an asymptotic coefficient associated with the Riemann
critical line or the distribution of zeta zeros.

A rigorous investigation of such an interpretation would
need to state the target asymptotic relation, derive it from
explicit assumptions, and compare it with known results or
independent numerical computations.

% ================================================================
% 7. RELATION TO THE TREE-SHADOW STUDY
% ================================================================
\section{Relation to the Tree-Shadow Study}

The companion study, \textit{The Tree and the Shadow},
considers the exact geometric identity
\begin{equation}
\tan\left(\frac{\pi}{8}\right)=\sqrt{2}-1
\approx0.41421356.
\label{eq:geometry}
\end{equation}

The identity in Equation~\eqref{eq:geometry} is exact.
Its application to tree crown morphology, however, is an
empirical hypothesis rather than a mathematical consequence
of trigonometry.

The companion study reports five mango-tree observations
with a sample mean crown ratio of approximately \(0.41297\)
and a TOST analysis using an equivalence margin of
\(\pm0.03\). Those results are exploratory because the sample
is small and comes from a single location.

The numerical proximity between the reported sample mean
and \(\sqrt{2}-1\) does not establish a causal relationship
between tree growth, solar geometry, and the Riemann zeta
function.

Similarly, the integer cutoff \(x_p=21\) is defined from
\(\lfloor\gamma_2\rfloor\). This definition does not, by
itself, establish a mathematical connection between
Equation~\eqref{eq:geometry} and the zeta zeros.

The proposed relationship among the geometric ratio,
biological observations, and numerical cutoff is therefore
a research hypothesis requiring independent tests.

% ================================================================
% 8. THE NORMALIZATION PROPOSAL
% ================================================================
\section{The Al-Hindara Normalization Proposal}

The project uses the relation
\begin{equation}
Z+C+A=1
\label{eq:unity}
\end{equation}
as a conceptual statement of unity.

For Equation~\eqref{eq:unity} to become a mathematical
identity, the three terms must be defined in a common
normalization scheme. The present paper does not supply a
derivation showing that the damping function, a separate
quantity \(C(x)\), and the accumulation function \(A(x)\)
sum to one.

In particular, setting
\[
C(x)=e^{-x/x_p}
\]
does not imply that
\[
Z(x)+C(x)+A(x)=1.
\]
That equality would have to be established independently,
either by construction with explicitly defined normalized
quantities or by a proof based on their definitions.

Accordingly, Equation~\eqref{eq:unity} remains a conceptual
framework in this paper, not a proven conservation law.

% ================================================================
% 9. DISCUSSION
% ================================================================
\section{Discussion: Evidential Status of the Claims}

The claims considered in this paper fall into three categories.

\subsection*{Exact mathematical results}

The definition \(x_p=\lfloor\gamma_2\rfloor=21\) follows from
the stated numerical value of \(\gamma_2\). The sign of
\(Z(x)\) on the intervals \((1,2)\) and \((2,\infty)\) follows
from the signs of the factors in Equation~\eqref{eq:field}.
The derivative relation \(A'(x)=Z(x)\) follows from the
Fundamental Theorem of Calculus.

\subsection*{Exploratory empirical observations}

The companion tree-shadow study reports a small pilot sample
whose mean is close to \(\sqrt{2}-1\). This observation is
not sufficient to establish a universal biological pattern.
Independent replication, transparent measurements, and a
suitable sampling design are necessary.

\subsection*{Conjectural modeling proposals}

The interpretation of \(x_p=21\) as a meaningful bridge
between geometry and number theory, the proposed coefficient
\(K\), and the normalization relation \(Z+C+A=1\) require
further work. No proof connecting the proposed damping
function to the Riemann zeta function is given here.

This separation is necessary to make the framework
scientifically testable and to avoid treating numerical
coincidence as proof of a mathematical relationship.

% ================================================================
% 10. CONCLUSION
% ================================================================
\section{Conclusion}

This paper defines a candidate damping function
\[
Z(x)=x^{-5}\ln(x)\sin\left(\frac{2\pi}{x}\right)e^{-x/21}
\]
using the integer \(21=\lfloor\gamma_2\rfloor\).

The function's sign can be determined analytically:
it is negative on \(1<x<2\) and positive on \(x>2\).
For the accumulation function
\[
A(x)=\int_1^x Z(t)\,dt,
\]
this means that \(A(x)\) has a minimum at \(x=2\) and is
increasing for all \(x>2\). Therefore, the stated model does
not support a maximum at \(x=42\).

The quantity \(K=2\pi e^{-1/42}\) is a well-defined
numerical expression, but its proposed interpretation
as a coefficient related to the Riemann critical line
remains unproved.

The companion tree-shadow study supplies an exploratory
empirical comparison with the exact value \(\sqrt{2}-1\).
That comparison does not establish a causal connection
between biological morphology and the zeta zeros.

Future work should focus on defining any alternative
accumulation function precisely, testing the model against
appropriate null models, providing reproducible numerical
code, and proving the normalization relation if it is
intended to be an exact mathematical statement.

% ================================================================
% REFERENCES
% ================================================================
\begin{thebibliography}{99}

\bibitem{edwards1974}
Edwards, H. M. (1974).
\textit{Riemann's Zeta Function}.
Academic Press.

\bibitem{schuirmann1987}
Schuirmann, D. J. (1987).
A comparison of the two one-sided tests procedure and the
power approach for assessing the equivalence of average
bioavailability.
\textit{Journal of Pharmacokinetics and Biopharmaceutics},
15(6), 657--680.

\bibitem{halle1978}
Hall\'e, F., Oldeman, R. A. A., and Tomlinson, P. B. (1978).
\textit{Tropical Trees and Forests: An Architectural Analysis}.
Springer-Verlag.

\bibitem{jabri2026frame2}
Al-Jabri, A. M. N. (2026).
\textit{The Tree and the Shadow: Geometric Derivation of the
Reciprocal Silver Ratio and an Exploratory Statistical Analysis
of Tree Crown Ratios}.
Companion manuscript. Publication status and persistent
identifier should be verified before formal submission.

\end{thebibliography}

% ================================================================
% APPENDIX
% ================================================================
\appendix

\section{Reproducibility Checklist}

For a reproducible computational study, archive the following:

\begin{enumerate}
\item The precise definition and domain of each function.
\item The numerical value used for \(\gamma_2\) and the
operation producing \(x_p=21\).
\item The integration method and numerical tolerance used
to evaluate \(A(x)\).
\item The search interval and grid resolution for any
reported numerical extremum.
\item Code that reproduces the function values, sign
analysis, and numerical integration.
\item A clear distinction between numerical observations
and analytically proved properties.
\item The definition and derivation of any proposed
normalization relation.
\end{enumerate}

\section{Supplementary Project Resources}

The following project resources were supplied by the author.
Their current availability and contents should be checked
before publication.

\begin{itemize}
\item \textbf{Zenodo Record (This Paper):} \\
\url{https://zenodo.org/records/23540283}

\item \textbf{DOI (This Paper):} \\
\url{https://doi.org/10.5281/zenodo.23540283}

\item Frame 2 visual simulation:
\url{https://jabri-com.vercel.app/tree/tree2.html}

\item Frame 2 TOST tool:
\url{https://jabri-com.vercel.app/tree/tree3.html}

\item Companion manuscript record (Frame 2):
\url{https://osf.io/yc9mb/}
\end{itemize}

\end{document}